\documentclass[../main.tex]{subfiles}

\begin{document}

\chapter{Evaluation}\label{ch:evaluation}

This chapter evaluates the proposed monitoring system. 
It presents the test setup, evaluates the correlation and detection capabilities, 
measures the performance overhead introduced by the interceptor, and concludes 
with a discussion of the threats to validity. 

\section{Test Setup}\label{sec:setup}

To facilitate integration with the Java testing framework, the entire test environment 
runs in containers. 
This comprises the Shibboleth IdP, an LDAP directory serving as the user store, the SPs, 
and the backend application of the monitoring tool.

Two SPs are deployed. 
The first is used for the detection scenarios and is configured to exhibit the weaknesses 
that the rules are intended to identify. 
The second is configured according to current recommendations and serves as a reference for 
the false positive rate, since any alert raised for a login against this provider is by 
definition a false positive. 

The evaluation distinguishes three levels of testing. 
Unit tests examine individual components in isolation, with their dependencies replaced by 
mocks, so that a test result reflects the behaviour of that component alone. 
Integration tests run against a real Elasticsearch instance, which additionally covers the 
persistence layer, the index mappings and the queries that a mock cannot represent. 
End to end tests exercise the complete pipeline, from a login at the SP through the log 
shipping stack to the alert produced by the backend. 

Logins are generated in two ways. 
An HTTP client harness performs the SAML login sequence directly, parsing the forms returned 
by the IdP and submitting them without rendering the pages. 
Selenium based tests drive an actual browser through the same sequence. 
The two differ in the interval between consecutive logins, since the harness proceeds as fast 
as the server responds whereas a browser adds the time required for page loading and form 
interaction. 
This difference is relevant for the correlation measurements in Section~\ref{sec:eval-correlation}, 
where the temporal density of logins determines how reliably records can be matched.

The environment is not a production deployment. 
It comprises a single IdP, two SPs and a small number of test accounts, and all traffic is generated 
for the purpose of measurement. 
The implications of this are discussed in Section~\ref{sec:validity}.

\section{Correlation}\label{sec:eval-correlation}

Section~\ref{subsec:impl-correlation} described the technical details of the correlation 
mechanism. 
This section presents the results of the correlation and evaluates the precision of the 
individual strategies.

To determine the precision of the less reliable strategies, a reference is required that 
indicates whether a given correlation was correct. 
The correlation via \texttt{INBOUND\_MESSAGE\_ID} establishes a one-to-one relationship, as 
explained in Section~\ref{subsec:impl-correlation}, since the identifier is unique per request. 
Wherever this strategy produced a match, the correct counterpart is therefore known without 
ambiguity, which allows it to serve as a ground truth for the remaining strategies.

The measurement method is as follows:

\begin{itemize}
    \item Collect all interceptor reports of a configurable time period for which 
          \texttt{correlated\_by = INBOUND\_MESSAGE\_ID}
    \item For each report and each window size, execute the exact query of the heuristic 
          strategy and select the same candidate as the productive system would, namely the 
          one with the smallest temporal distance
    \item Compare the selected entry against the known correct audit identifier, classifying 
          the outcome as a match, an incorrect assignment, or no match
    \item Additionally record the temporal offset between the reception of the interceptor 
          report and the actual audit event time as a distribution
\end{itemize}

The underlying data source is 200 consecutive logins of a single test user via the HTTP login 
harness, preceded by three warm-up logins that are included in the measurement. 
In terms of candidate ambiguity this represents the worst case, since all candidates share the 
same \textit{subjectHash}, \textit{ipHash} and SP, leaving time as the only distinguishing 
feature.

\subsection{Defects Identified during the Evaluation}\label{subsec:eval-correlation-findings}

Before the evaluation, 3{,}690 out of 3{,}690 correlations had been established through the 
hashed identity and time window strategy, and not a single one through 
\texttt{INBOUND\_MESSAGE\_ID}. 
The cause was not located in the application but in the audit format of the IdP. 
The \textit{shibboleth.AuditFormattingMap} in \textit{conf/audit.xml} did not contain 
\texttt{\%I}. 
The Logstash pipeline was already mapping \textit{[shib][I]} to the corresponding field, but it 
was waiting for a value that the configuration never produced. 
The precondition of the strategy, namely the presence of the inbound message identifier in the 
field registry, therefore remained permanently unmet. 
After adding \texttt{\%I} to the audit format, the registry recognised the field and the 
strategy activated itself during operation, without restarting the system. 
This is direct evidence of the dynamic field adaptation described in 
Section~\ref{sec:impl-aggregation}.

A second finding concerns the timestamp used by the time window strategies. 
These compared against the \textit{timestamp} attribute, which reflects the ingest time assigned 
by Filebeat rather than the time of the authentication event. 
Since Filebeat processes log lines in batches, this value is shared by all entries of a batch. 
A measurement on benchmark traffic showed that 2{,}423 audit entries were distributed across only 
448 distinct \textit{timestamp} values, an average of 5.4 events per timestamp. 
Within such a group the temporal distance is identical for all candidates, so the chronological 
order of several logins could not be reconstructed.

This was resolved by introducing a dedicated field \textit{event\_time}, which Logstash derives 
from \texttt{\%T}, the event time recorded by the IdP, using a date filter. 
Both time window strategies have operated on this field since. 
The \textit{timestamp} attribute remains the basis for scheduler markers and the indexing grace 
period, where the ingest time is the factually correct reference.

\subsection{Results of the Evaluation}\label{subsec:eval-correlation-results}

With \textit{event\_time} as the basis, the temporal resolution improves to 1.0 events per 
timestamp. 
Each event therefore carries a distinct value, in contrast to the 5.4 events per timestamp 
observed with the ingest time.
Table~\ref{tab:timestamp-precision} and Table~\ref{tab:event-time-precision} show the precision 
of the heuristic strategy against the exact ground truth before and after this change. 
The two are based on separate runs, since the change to \textit{event\_time} required the data 
to be collected anew. 

\begin{table}[h]
\centering
\caption{Precision of \texttt{HASH\_AND\_TIME\_WINDOW} based on the Filebeat ingest time, single user via HTTP harness (n = 415)}\label{tab:timestamp-precision}
\small
\begin{tabular}{lrrrrrr}
\toprule
Window & Hits & Mismatch & No Match & Precision & Coverage & Candidates \\
\midrule
$\pm$2\,s  & 24 & 373 & 18 & 6,05\% & 95,66\% & 11,79  \\
$\pm$5\,s  & 27 & 384 & 4  & 6,57\% & 99,04\% & 28,90  \\
$\pm$10\,s & 29 & 386 & 0  & 6,99\% & 100\%   & 55,47  \\
$\pm$30\,s & 29 & 386 & 0  & 6,99\% & 100\%   & 138,48 \\
$\pm$60\,s & 29 & 386 & 0  & 6,99\% & 100\%   & 192,60 \\
\bottomrule
\end{tabular}
\end{table}

\begin{table}[h]
\centering
\caption{Precision of \texttt{HASH\_AND\_TIME\_WINDOW} based on the IdP event time, single user via HTTP harness (n = 203)}\label{tab:event-time-precision}
\small
\begin{tabular}{lrrrrrr}
\toprule
Window & Hits & Mismatch & No Match & Precision & Coverage & Candidates \\
\midrule
$\pm$1\,s  & 203 & 0 & 0 & 100\% & 100\% & 6,42   \\
$\pm$2\,s  & 203 & 0 & 0 & 100\% & 100\% & 12,36  \\
$\pm$5\,s  & 203 & 0 & 0 & 100\% & 100\% & 30,24  \\
$\pm$10\,s & 203 & 0 & 0 & 100\% & 100\% & 58,12  \\
$\pm$30\,s & 203 & 0 & 0 & 100\% & 100\% & 144,68 \\
$\pm$60\,s & 203 & 0 & 0 & 100\% & 100\% & 197,60 \\
\bottomrule
\end{tabular}
\end{table}

\begin{table}[h]
\centering
\caption{Temporal offset between interceptor report reception and IdP event time (n = 203)}\label{tab:time-offset}
\small
\begin{tabular}{lllll}
\toprule
min & p50 & p95 & max & Span \\
\midrule
$-98$\,ms & $-32$\,ms & $-23$\,ms & $+75$\,ms & 173\,ms \\
\bottomrule
\end{tabular}
\end{table}

Table~\ref{tab:time-offset} shows the distribution of this offset. 
The predominantly negative offset is consistent with the position of the interceptor in the 
profile flow. 
It runs as a \textit{postAuthenticationFlow} in the middle of the profile, whereas \texttt{\%T} 
marks the completion of the transaction. 
The report therefore typically reaches the backend around 30\,ms before the IdP writes the 
corresponding audit entry.

\subsubsection{Precision of TIME\_WINDOW under True Candidate Ambiguity}\label{subsubsec:precision-time-window-multi-user}

The measurements in Table~\ref{tab:event-time-precision} are based on a single test user. 
In that case the candidate sets of \texttt{HASH\_AND\_TIME\_WINDOW} and \texttt{TIME\_WINDOW} 
are identical, since a second \textit{usernameHash} never occurs within the observation period. 
The weaker strategy is therefore never actually tested, only carried along. 
To address this, the same measurement methodology was applied to a dataset with several test 
users active in parallel.

The extended dataset comprises three users with 200 logins each, started simultaneously against 
the same SP, resulting in 600 reports in total. 
The ground truth was cleared before the run to prevent any mixing with the single user data. 
The results are summarised in Table~\ref{tab:precision-strategies}.

\begin{table}[h]
\centering
\caption{Precision of both strategies under three concurrent users via the HTTP harness (n = 600)}\label{tab:precision-strategies}
\small
\begin{tabular}{lrrr}
\toprule
Strategy & Precision (all windows 1--60\,s) & Coverage & Candidates ($\pm$1\,s) \\
\midrule
\texttt{HASH\_AND\_TIME\_WINDOW} & 100\% & 100\% & 5,66  \\
\texttt{TIME\_WINDOW}            & 65\%  & 100\% & 16,26 \\
\bottomrule
\end{tabular}
\end{table}

The precision of \texttt{TIME\_WINDOW} is identical across all six measured window sizes at 
65\%. 
This follows directly from the selection rule, according to which the candidate with the 
smallest temporal distance is chosen. 
Once the correct candidate lies within the window, a more distant candidate admitted by a larger 
window can no longer displace it. 
Since the number of reports without a match is consistently zero, the $\pm$1\,s window already 
sufficed to include the correct candidate, only not always as the closest one.

\begin{table}[h]
\centering
\caption{Precision of \texttt{TIME\_WINDOW} by local login density, HTTP harness, window = 10\,s}\label{tab:local-density}
\small
\begin{tabular}{lrrrr}
\toprule
Density Bucket & n & Hits & Mismatch & Precision \\
\midrule
< 100\,ms      & 521 & 311 & 210 & 59,69\% \\
100--300\,ms   & 65  & 65  & 0   & 100\%   \\
300\,ms--1\,s  & 14  & 14  & 0   & 100\%   \\
> 1\,s         & 0   & 0   & 0   & ---     \\
\bottomrule
\end{tabular}
\end{table}

Table~\ref{tab:local-density} shows that the incorrect assignments of \texttt{TIME\_WINDOW} are 
concentrated entirely among logins that lie less than 100\,ms apart from the next login of the 
same SP. 
521 of the 600 reports fall into this category. 
Above an interval of 100\,ms the strategy is error-free within the measured dataset.

\subsubsection{Precision under Real Browser Traffic}\label{subsubsec:precision-browser-traffic}

All measurements presented so far are based on the synthetic HTTP login harness. 
To examine whether its timing behaviour is representative of real login processes, the same 
measurement methodology was applied to a second dataset generated through Chrome-based logins 
using Testcontainers and Selenium instead of a raw HTTP client. 
The data source is 202 logins of a single test user. 
The results are summarised in Table~\ref{tab:e2e-single-user-results}, where the HTTP client 
column repeats the values of Table~\ref{tab:event-time-precision}.

\begin{table}[h]
\centering
\caption{Single user, HTTP harness compared to real browser logins}\label{tab:e2e-single-user-results}
\small
\begin{tabular}{lrr}
\toprule
  & HTTP client (n = 203) & Browser (n = 202) \\
\midrule
Precision (both strategies) & 100\%   & 100\%  \\
Candidates at $\pm$1\,s     & 6,42    & 1,00   \\
Candidates at $\pm$60\,s    & 197,60  & 42,61  \\
\bottomrule
\end{tabular}
\end{table}

\begin{table}[h]
\centering
\caption{Local login density under real browser traffic, single user (n = 202)}\label{tab:e2e-local-density}
\small
\begin{tabular}{lrr}
\toprule
Density Bucket & n & Precision \\
\midrule
< 100\,ms      & 0   & ---   \\
100--300\,ms   & 0   & ---   \\
300\,ms--1\,s  & 0   & ---   \\
> 1\,s         & 202 & 100\% \\
\bottomrule
\end{tabular}
\end{table}

As Table~\ref{tab:e2e-local-density} shows, all 202 logins lie more than one second apart from 
their nearest neighbour. 
All 202 logins lie more than one second apart from their nearest neighbour. 
The candidate mean at $\pm$1\,s is 1.00 instead of 6.42, which places the login density roughly 
six times lower than with the HTTP client harness. 
This is consistent with the additional time required for real browser interaction, comprising 
page loading, form handling and the overhead introduced by Selenium, compared to a script 
without waiting times. 
At this density \texttt{HASH\_AND\_TIME\_WINDOW} and \texttt{TIME\_WINDOW} yield identical 
results, since with an average of one candidate per window the additional hash filter has 
nothing left to exclude, for the same reason as in the single user case.

The offset between report reception and event time ranges from $-102$ to $-17$\,ms (n = 202), 
which is narrower and consistently negative compared to the range reported in 
Table~\ref{tab:time-offset}. 
With a single run of this size no reliable trend can be established, but there is equally no 
indication that real logins are temporally more irregular than the synthetic client.

\subsubsection{Precision among Multiple Simultaneous Browser Users}\label{subsubsec:precision-multiple-browser-users}

What remains open is the combination of the two preceding measurements, namely several users and 
real browser traffic at the same time. 
The same measurement methodology was therefore applied with three test users running 
simultaneously against the same SP. 
The data source is 202 logins each, comprising 200 measured logins and the hard-coded warm-up of 
two, resulting in 606 reports in total. 
The results are summarised in Table~\ref{tab:precision-coverage-strategies}, 
Table~\ref{tab:sum-local-density} and Table~\ref{tab:comparison-multi-user}.

\begin{table}[h]
\centering
\caption{Precision and coverage under three concurrent browser users (n = 606)}\label{tab:precision-coverage-strategies}
\small
\begin{tabular}{lrr}
\toprule
Strategy & Precision (all windows 1--60\,s) & Coverage \\
\midrule
\texttt{HASH\_AND\_TIME\_WINDOW} & 100\%   & 100\% \\
\texttt{TIME\_WINDOW}            & 89,93\% & 100\% \\
\bottomrule
\end{tabular}
\end{table}

\begin{table}[h]
\centering
\caption{	Precision by local login density under three concurrent browser users, window = 10\,s}\label{tab:sum-local-density}
\small
\begin{tabular}{lrrr}
\toprule
Density Bucket & n & \texttt{HASH\_AND\_TIME\_WINDOW} & \texttt{TIME\_WINDOW} \\
\midrule
< 100\,ms      & 138 & 100\% & 55,80\% \\
100--300\,ms   & 132 & 100\% & 100\%   \\
300\,ms--1\,s  & 258 & 100\% & 100\%   \\
> 1\,s         & 78  & 100\% & 100\%   \\
\bottomrule
\end{tabular}
\end{table}

\begin{table}[h]
\centering
\caption{Comparison with the synthetic multi-user dataset}\label{tab:comparison-multi-user}
\small
\begin{tabular}{lrr}
\toprule
  & Synthetic (n = 600) & Real browser (n = 606) \\
\midrule
Precision in the < 100\,ms bucket          & 59,69\%          & 55,80\%          \\
Share of logins below 100\,ms              & 86,8\% (521/600) & 22,8\% (138/606) \\
Overall precision of \texttt{TIME\_WINDOW} & 65,00\%          & 89,93\%          \\
\bottomrule
\end{tabular}
\end{table}

The precision within the critical sub-100\,ms zone is nearly identical for real and synthetic 
traffic, at approximately 56 to 60 percent. 
The weakness measured in Section~\ref{subsubsec:precision-time-window-multi-user} is therefore 
not an artifact of the HTTP client timing, but a reproducible density effect. 
What differs is how frequently this zone is reached at all. 
The artificially dense stress test enters it in 87\% of cases, whereas three genuinely concurrent 
browser users do so in only 23\%. 
The lower overall precision of the synthetic dataset therefore describes the stress case 
correctly, but overstates how often it occurs under the observed user load.

\subsubsection{The Independent Contribution of \textit{ipHash}}\label{subsubsec:influence-of-iphash}

It remained unclear whether \textit{ipHash} actually contributed to the separation in the 
multi-user datasets, or whether \textit{usernameHash} alone accounted for the measured precision. 
The initial assumption had been that all benchmark processes would share a common 
\textit{ipHash}. 
A control query on the raw data refuted this for the browser path, since each 
\texttt{BrowserWebDriverContainer} started in parallel is assigned its own Docker bridge 
address. 
The three test users in Section~\ref{subsubsec:precision-multiple-browser-users} therefore had 
three distinct \textit{ipHash} values rather than a shared one.

This could be verified against the existing ground truth without performing a new login run. 
\texttt{CorrelationPrecisionAnalysis} was extended by a third, purely analytical strategy that 
searches by SP, \textit{ipHash} and time window while omitting \textit{usernameHash}. 
This is not a productive strategy but a diagnostic variant introduced to isolate the 
contribution of \textit{ipHash}. 
The data source is the same 606 reports as in 
Section~\ref{subsubsec:precision-multiple-browser-users}, reanalysed with the additional 
strategy. 
Table~\ref{tab:ip-hash-correlation} shows the result of this comparison.

\begin{table}[h]
\centering
\caption{Precision with and without \textit{usernameHash} (n = 606)}\label{tab:ip-hash-correlation}
\small
\begin{tabular}{lrr}
\toprule
Strategy & Precision (all windows) & Candidates ($\pm$1\,s) \\
\midrule
\texttt{HASH\_AND\_TIME\_WINDOW} & 100\% & 1,00 \\
\texttt{IP\_ONLY}                & 100\% & 1,00 \\
\bottomrule
\end{tabular}
\end{table}

The comparison across all six window sizes shows not merely a convergence but exact agreement 
down to the second decimal place, as Table~\ref{tab:ip-hash-candidates} shows.

\begin{table}[h]
\centering
\caption{Mean candidate count per window (n = 606)}\label{tab:ip-hash-candidates}
\small
\begin{tabular}{lrr}
\toprule
Window & \texttt{HASH\_AND\_TIME\_WINDOW} & \texttt{IP\_ONLY} \\
\midrule
$\pm$1\,s  & 1,00  & 1,00  \\
$\pm$2\,s  & 2,19  & 2,19  \\
$\pm$5\,s  & 3,86  & 3,86  \\
$\pm$10\,s & 7,49  & 7,49  \\
$\pm$30\,s & 21,67 & 21,67 \\
$\pm$60\,s & 41,58 & 41,58 \\
\bottomrule
\end{tabular}
\end{table}

Both strategies also coincide when broken down by local login density, including in the 
otherwise critical bucket below 100\,ms, where both achieve 138 of 138 matches. 
In this dataset \textit{usernameHash} never once excludes a candidate that \textit{ipHash} had 
not already ruled out on its own, because every browser container had a stable address of its 
own and each \textit{ipHash} bucket was therefore already unambiguous.

This shows that \textit{ipHash} separates distinct sources, whether different devices, networks 
or, as here, containers, just as effectively as the full combination of both identifiers. 
It does not show what happens when several users share one address, as is the case behind 
corporate NAT or on a public wireless network. 
That scenario could not be produced with the existing test setup, which assigns one container 
per user, and is examined in the following section.

\subsubsection{Precision among Multiple Users Sharing One Address}\label{subsubsec:precision-nat}

The preceding section established the contribution of \textit{ipHash} only for separate sources. 
The opposite case, several users behind a shared address, is at least as relevant in practice 
and could not be reproduced with the container based setup. 
A dedicated measurement tool was therefore developed that runs multiple users as parallel 
threads within a single process rather than in separate containers. 
Since all threads send from the same machine, they necessarily share a source address, which 
follows from the architecture rather than from any network configuration.

Before the evaluation, a cardinality aggregation on \textit{ip\_hash} over the period of the run 
confirmed exactly one distinct value across all 600 documents. 
The shared address is therefore measured rather than assumed. 
The data source is three test users with 200 logins each, started from the same process. 
The results are summarised in Table~\ref{tab:precision-strategy-nat} and 
Table~\ref{tab:precision-strategy-nat-local-density}. 

\begin{table}[h]
\centering
\caption{Precision of both strategies with three users sharing one source address (n = 600)}\label{tab:precision-strategy-nat}
\small
\begin{tabular}{lrr}
\toprule
Strategy & Precision (all windows 1--60\,s) & Coverage \\
\midrule
\texttt{HASH\_AND\_TIME\_WINDOW} & 100\%   & 100\% \\
\texttt{TIME\_WINDOW}            & 73,67\% & 100\% \\
\bottomrule
\end{tabular}
\end{table}

\begin{table}[h]
\centering
\caption{Precision by local login density with a shared source address, window = 10\,s}\label{tab:precision-strategy-nat-local-density}
\small
\begin{tabular}{lrrr}
\toprule
Density Bucket & n & \texttt{HASH\_AND\_TIME\_WINDOW} & \texttt{TIME\_WINDOW} \\
\midrule
< 100\,ms      & 504 & 100\% & 68,65\% \\
100--300\,ms   & 96  & 100\% & 100\%   \\
300\,ms--1\,s  & 0   & ---   & ---     \\
> 1\,s         & 0   & ---   & ---     \\
\bottomrule
\end{tabular}
\end{table}

\texttt{HASH\_AND\_TIME\_WINDOW} remains at 100\% despite the shared address, since 
\textit{usernameHash} carries the entire separation on its own, irrespective of whether 
\textit{ipHash} contributes anything.

The diagnostic strategy of the preceding section serves as a cross-check and yields exactly the 
same figures here as \texttt{TIME\_WINDOW}, identical for every window size and every density 
bucket at 73,67\% precision with 442 matches and 158 incorrect assignments. 
This is not a coincidence but the necessary consequence of a single shared address. 
A filter for the only \textit{ipHash} value occurring in the entire period excludes no candidate 
that would not already fall within the time window. 
In this dataset \textit{ipHash} therefore demonstrably contributes nothing to the separation, 
which is the exact counterpart to the preceding section, where \texttt{IP\_ONLY} coincided with 
\texttt{HASH\_AND\_TIME\_WINDOW} because the address alone already sufficed.

The precision of \texttt{TIME\_WINDOW} at 73,67\% falls between the 65,00\% measured in 
Section~\ref{subsubsec:precision-time-window-multi-user} and the 89,93\% measured in 
Section~\ref{subsubsec:precision-multiple-browser-users}, which is consistent with the observed 
density. 
504 of 600 logins, or 84\%, lie less than 100\,ms apart, close to the 86,8\% of the synthetic 
harness and well above the 22,8\% of the browser based measurement. 
Several threads running without waiting times from a single process therefore generate a density 
comparable to the serial HTTP harness rather than to real browser interaction. 
The effect stems primarily from login density and not from the shared address itself, which is 
consistent with the conclusion already drawn in 
Section~\ref{subsubsec:precision-multiple-browser-users}.
The offset between report reception and event time ranges from $-50$ to $+110$\,ms with a median 
of $-26$\,ms (n = 600), which is of the same order of magnitude as in the other measurements.

\subsection{Interpretation of the Results}\label{subsec:interpretation-correlation}

Window size is not the limiting factor. Precision remains at 100\% even though the mean 
number of candidates rises from 6.4 at $\pm$1,s to 197.6 at $\pm$60,s.
Before turning to precision, the coverage of the correlation is considered, that is the 
proportion of interceptor reports for which a counterpart was found at all. 
Across the measured period, 3{,}762 of 3{,}765 reports were correlated, leaving three that 
remained without a counterpart after the retry window had elapsed. 
The value is only determinable because the sweep described in 
Section~\ref{subsec:impl-correlation} marks expired reports as permanently uncorrelated. 
Without that distinction, a report still awaiting its audit entry would be indistinguishable 
from one that will never receive it, and the denominator of the ratio would change with every 
query rather than remaining fixed.

Coverage is also independent of the strategy that produced a match.  
In every measurement reported above it remains at 100\% from a window of $\pm$5\,s onwards, 
including for \texttt{TIME\_WINDOW}, which loses precision but not coverage under high login 
density. 
The two properties therefore fail independently of one another. 
A weaker strategy does not leave reports unassigned, it assigns them to the wrong counterpart.

Window size is not the limiting factor. 
Precision remains at 100\% even though the mean number of candidates rises from 6.4 at 
$\pm$1\,s to 197.6 at $\pm$60\,s. 
What matters is solely that the temporal jitter is smaller than the interval between two 
consecutive logins of the same combination of subject, address and SP.

This yields a quantifiable operational limit. 
With a jitter span of 173\,ms and an observed login interval of approximately 330\,ms, the limit 
lies at about six logins per second per tuple of subject, address and SP. 
Beyond this point the heuristic can no longer resolve the order reliably, whereas the exact 
strategy is unaffected by it.

The measurements also support a concrete configuration recommendation, namely 
\textit{lg.correlation.window-seconds=1}. 
Coverage and precision remain at 100\%, but 6.4 candidates are examined per report instead of 
58.1, which reduces the effort by a factor of nine on the basis of evidence rather than a 
fixed assumption.

The fallback chain has a quantifiable value. 
Where the unique identifier is absent from the audit log, for instance because it is not recorded 
for data protection reasons, the heuristic delivers equivalent results under the measured 
conditions. 
The loss of information is therefore not qualitative but quantitative, and it only takes effect 
above the login density stated above.

The cost of omitting the hashes is thus quantified rather than merely postulated. 
This applies, however, only to the comparison between the exact strategy and 
\texttt{HASH\_AND\_TIME\_WINDOW}. 
If both \textit{ipHash} and \textit{usernameHash} are missing, so that \texttt{TIME\_WINDOW} has 
to step in, the picture changes. 
Section~\ref{subsubsec:precision-time-window-multi-user} shows a drop in precision from 100\% to 
65\% with three simultaneously active users, concentrated on logins less than 100\,ms apart. 
For a SP with several concurrent users, \texttt{TIME\_WINDOW} is therefore to be classified as a 
stopgap for the case of missing hashes rather than as an equivalent alternative, unlike what the 
single user measurement in Section~\ref{subsec:eval-correlation-results} alone would suggest, 
where both strategies are indistinguishable.

There is a first indication that the operational limit is not reached in practice. 
Section~\ref{subsubsec:precision-browser-traffic} measures the login density under real browser 
traffic instead of the synthetic HTTP client harness, and all 202 logins lie more than one second 
apart, well below the limit of roughly six logins per second calculated above. 
For a single automated browser user the synthetic worst case of 
Section~\ref{subsec:eval-correlation-results} is therefore considerably denser than real timing, 
although this has not yet been shown for several concurrent real users.

What matters is the distribution of density, not merely the possibility of a collision. 
Section~\ref{subsubsec:precision-multiple-browser-users} shows that the weakness of 
\texttt{TIME\_WINDOW} in the sub-100\,ms zone measured in 
Section~\ref{subsubsec:precision-time-window-multi-user} is just as pronounced under real traffic, 
at approximately 56\% precision, and is therefore not a measurement artifact. 
What differs is how often that zone is reached. 
Real concurrent users produce it noticeably less frequently, in 23\% of logins, than the 
artificially dense stress case, which reaches it in 87\%. 
For operation this means that the risk is real and reproducibly demonstrated, but occurs less 
often under the observed user load than the stress case figure alone would suggest.

The contribution of \textit{ipHash} depends entirely on whether the source actually separates 
users. 
Section~\ref{subsubsec:influence-of-iphash} shows that with separate sources \textit{ipHash} alone 
separates just as well as \textit{ipHash} and \textit{usernameHash} combined, yielding identical 
candidate sets in every measured window. 
Section~\ref{subsubsec:precision-nat} shows the counterpart, namely that where several users share 
one address \textit{ipHash} contributes nothing at all, since \texttt{IP\_ONLY} and 
\texttt{TIME\_WINDOW} produce identical figures there. 
For scenarios in which \textit{usernameHash} is unavailable, this is the differentiated answer. 
The heuristic is only as good as the least discriminating identifier available to it. 
Where a shared address and a missing user identifier coincide, only \texttt{TIME\_WINDOW} remains, 
with the density dependence observed consistently throughout.

The productive strategy withstands the shared address case as long as \textit{usernameHash} 
remains available. 
This is the practically most relevant conclusion of Section~\ref{subsubsec:precision-nat}, since 
\texttt{HASH\_AND\_TIME\_WINDOW} retains 100\% precision although all three test users share one 
address. 
Only when the user identifier is absent in addition to the shared address does the correlation 
become a question of the time window alone, which reduces to the density dependence already 
established rather than constituting a new risk specific to network address translation.

\section{Detection}\label{sec:eval-detection}

The three test layers and what each can check
Coverage: thirteen out of thirteen on layer one, five on layer two
Why not all of them on layer two
The Logstash error and the fact that only layer two could detect it
The coverage analysis against the FIM survey, including the unaddressed lines
False-positive baseline, as soon as it's measured
Honestly: What these numbers prove—and what they don't

\section{Performance Overhead}\label{sec:overhead}

This section presents the results of the overhead evaluation of the custom interceptor. 
The interceptor runs as a \textit{postAuthenticationFlow} within the IdP and therefore 
lies on the critical path of each authentication process.
The evaluation addresses two questions:
\begin{itemize}
      \item How much overhead does the interceptor introduce to the end-to-end authentication latency?
      \item Does the difference between asynchronous and synchronous transmission justify the use of fire-and-forget transmission?
\end{itemize}

To measure the overhead introduced by the interceptor, three different configuration modes 
are evaluated: \textit{baseline}, \textit{async}, and \textit{sync}.
In the \textit{baseline} configuration, no \textit{postAuthenticationFlow} is configured for the SP. 
The custom interceptor is therefore not executed during the authentication process.
In the \textit{async} configuration, the \textit{postAuthenticationFlow} is enabled for the SP and 
the custom interceptor is executed. 
The interceptor transmits the collected data to the monitoring tool without waiting for a response. 
The transmission therefore follows a fire-and-forget approach.
The \textit{sync} configuration also executes the custom interceptor for the SP, but the interceptor 
waits for a response from the monitoring tool before completing its execution. 
It therefore verifies whether the transmitted payload was received successfully.

The overhead is evaluated using three measurements.
First, an \textit{internal} measurement is obtained through the \textit{executionTimeMs} instrumentation 
within the interceptor. 
It measures the time required for data collection, hashing, and XML extraction performed by the action. 
This represents a lower-bound measurement because payload construction and data transmission are excluded. 
The measured value is included in the generated payload.
Second, an \textit{external} benchmark is performed by the test class and measures the complete SAML 
authentication process over HTTP. 
This measurement includes the effects of the interceptor and the surrounding system and therefore provides 
an upper-bound estimate of the introduced overhead, although with lower resolution.
The two measurements complement each other. 
The internal measurement provides a precise measurement of the interceptor's execution time, while the 
external measurement captures the overall impact on the authentication process.
Finally, CPU utilisation is measured to evaluate the additional computational load introduced by the 
interceptor. 

A preliminary run showed three different settings within a single long run. 
Iteration 0 to 80, showed a JIT warum up section with falling latency. 
Iteration 80 to 285, had a stable plateau. 
All subsequent iterations showed a drop in latency. 
A residual increase persisted starting around iteration 125, which is why the measurement window was 
chosen conservatively. 
A warum with 100 iterations and the actual calculation with 100 iterations. 

To eliminate position biasy as best as possible a rotating configuration order was implemented. 
Six runs with a rotating order, so that each configuration occurs exactly twice in each measurement 
position. 
Balancing was necessary because a position drift occurs within a session. 
The experimental unit is the run, not the individual login. 
The 100 logins within a run are pseudo-replicates, and the dominant source of variation is the variation 
between runs. 


The Three Measurement Sources: Internal, External, CPU
The Evaluation Unit Is the Run, Not the Login—and Why
The Block Design to Counter Position Drift
Results for Async and Sync, with Confidence Intervals
CPU Results as Independent Confirmation
The Two Errors in My Own Analysis Script

\section{Threats to Validity}\label{sec:validity}

n gleich sechs, ein einzelner Lauf je Messung
Synthetischer Verkehr, ein Nutzer, ein SP bei den meisten Messungen
Kein gelabelter Produktivdatensatz, daher keine Erkennungsgüte bestimmbar
Die Zirkularität: Regeln und Testfälle aus derselben Hand -> nach prof angeblich nicht so schlimm
Schwellenwerte gesetzt, nicht empirisch bestimmt

\end{document}